%!TEX root = ../main.tex

\chapter{Prompt per Prism}\label{ch:prompt-prism}

Questa sezione contiene un prompt operativo da usare in Prism per trasformare il dossier preparatorio in una prima bozza della relazione finale. Prima dell'uso e opportuno completare i dati amministrativi indicati in appendice e sostituire i segnaposto presenti nel frontespizio.

\section{Prompt consigliato}

Scrivi una relazione finale di tirocinio in italiano, con tono universitario, chiaro e professionale, relativa al progetto \builder. Usa come base le informazioni raccolte in questo dossier e organizza il testo in circa dieci pagine di contenuto principale.

La relazione deve presentare \builder{} come applicazione web per la progettazione, visualizzazione, traduzione ed esportazione di schemi Entity-Relationship e schemi logici. Descrivi il contesto del tirocinio, gli obiettivi iniziali, le funzionalita realizzate o consolidate, le tecnologie utilizzate, le attivita di verifica e le competenze acquisite.

Valorizza il ruolo della Prof.ssa Alessandra Lumini come tutor del tirocinio, evitando citazioni dirette non documentate. Spiega che la supervisione ha contribuito a orientare le priorita del lavoro, con attenzione alla chiarezza dell'applicazione, alla qualita dell'interfaccia, all'usabilita e alla correttezza concettuale degli schemi prodotti.

Mantieni un equilibrio tra descrizione tecnica e riflessione formativa. Non trasformare il testo in una guida utente dettagliata: usa le funzionalita del progetto per spiegare il percorso svolto, le decisioni progettuali, le difficolta incontrate, i risultati raggiunti, i limiti ancora presenti e i possibili sviluppi futuri.

\section{Vincoli di stesura}

\begin{itemize}
    \item Usa un lessico coerente con una relazione di tirocinio universitaria.
    \item Non inventare date, ore, CFU, dati dello studente o struttura ospitante: lascia segnaposto dove le informazioni non sono disponibili.
    \item Non attribuire alla tutor frasi o giudizi diretti non documentati.
    \item Inserisci riferimenti prudenti a test, export, salvataggio, localizzazione, reverse engineering SQL e trasformazione verso lo schema logico.
    \item Concludi con una valutazione critica che distingua risultati ottenuti, limiti e sviluppi futuri.
\end{itemize}
